\section{Conclusion}

Ordering the period-six chiral H\'enon doublet repairs the loss of point
chronology.  A directed edge reconstructs the radical, the six-cycle, and
the full cubic splitting field.  The resulting curve is geometrically
connected, has deck group $\mathsf D_6$ of order twelve, and has genus one.
Yet its exact order-six time generator is an unramified torsion translation
and acts identically on $H^1$.  Positive genus and genuine point dynamics do
not by themselves produce a nontrivial cohomological time spectrum.

The cross-period comparison sharpens the same lesson.  In the explicitly
certified chiral sample through period seven, the first witnessed nontrivial
weight-one time sector occurs on the adopted period-seven component, not on
the unique period-six chiral doublet.  A simultaneous quadratic marker
signal does not repair the gap: both marker curves are affine copies of a
period-one fixed-point curve.  Their common field is inherited, not a
primitive clock bridge.

These are useful negative results for a Hilbert--P\'olya search.  They give
two early rejection tests: compute the actual time representation on
cohomology before interpreting a periodic cover spectrally, and factor any
cross-period marker coincidence through lower-period loci before treating
it as arithmetic structure.  C21 stops at those tests.  It supplies neither
a repetition law nor a determinant, and the next round should move to a
different dynamical form if no connected varying-period tower can be found.
